%% GENERATED by tools/gen_latex.py -- edit content/agents10.py instead.
%% Build:  latexmk -pdf agents10.tex     (or: pdflatex agents10.tex, twice)
\documentclass[aspectratio=169,11pt,t]{beamer}
\usepackage{itbpro}

\coursetitle{Agen kode}
\coursesubtitle{Kode adalah alat yang bisa menjadi alat apa pun --- itu sebabnya ia paling berguna, dan itu sebabnya ia satu-satunya yang butuh dinding sungguhan.}
\coursekicker{Chapter 10}
\coursesource{Grootendorst \& Alammar, An Illustrated Guide to AI Agents --- bab 10}
\courseduration{3 jam (2 sesi)}
\coursepresenter{Hendri Karisma}{Instructor}
\coursebrand{AI for Professional \textperiodcentered\ ITB \textperiodcentered\ Ch 10}

\title{Agen kode}
\author{Hendri Karisma}
\date{}

\begin{document}
\covertitle

\framekicker{Learning outcomes}
\begin{frame}[shrink=25]{Session Objectives}
\leadpar{By the end of this session, participants can:}
\begin{isteps}
\item \textbf{Menjelaskan kenapa menulis kode mengubah sifat jawaban} --- dari ditebak jadi dihitung, dan dari tak terbaca jadi bisa diperiksa.
\item \textbf{Menyebutkan lima jenis alat kode} dan memisahkan yang baca dari yang menulis.
\item \textbf{Menghitung kenapa seluruh repo tidak muat}, dan menjelaskan pola peta-cari-baca.
\item \textbf{Menyebutkan lapis sandbox} beserta apa yang TIDAK ditutupnya.
\item \textbf{Membandingkan alur tetap dengan agen bebas} untuk tugas rekayasa perangkat lunak.
\item \textbf{Menyebutkan risiko khas} agen kode yang tidak muncul pada agen lain.
\end{isteps}

\vskip 6pt
\vskip 4pt
\reslink{Site}{Course home}{https://hendrikarisma.my.id/rs/ai-products-course/index.html}
\reslink{Repo}{ai-agentic-demo --- kasus code\_analysis}{https://github.com/situkangsayur/ai-agentic-demo}
\reslink{Book}{Maarten Grootendorst \& Jay Alammar, An Illustrated Guide to AI Agents (O'Reilly Media, Early Release)}{https://www.oreilly.com/library/view/an-illustrated-guide/9798341662681/}
\vskip 2pt
\end{frame}

\coursesection{01}{Kenapa kode}{Bukan karena penggunanya programmer.}

\framekicker{Dasar}
\begin{frame}[shrink=25]{Menulis kode mengubah sifat jawabannya}
\begin{center}
\includegraphics[width=0.92\linewidth,height=0.60\textheight,keepaspectratio]{figs/agents10-whycode.tex.pdf}
\end{center}
\vskip 3pt{\color{ink3}\fontsize{7.4}{9.4}\selectfont Jawaban yang dihitung, dengan cara menghitung yang bisa dibaca ulang.\par}
{\fontsize{9.4}{12.6}\selectfont\color{fg2}Ketika model menjawab pertanyaan data secara langsung, jawabannya adalah \textbf{tebakan yang masuk akal}. Ketika ia menulis kode yang lalu dijalankan, jawabannya \textbf{hasil perhitungan} --- dan cara menghitungnya ada di layar untuk diperiksa.\par}\vskip 4pt
\begin{iband}[signal]
Ini alasan yang sama kenapa modul ini terus mengulang \textit{serahkan aritmetika ke kode}. Bedanya di sini: \textbf{kodenya ditulis saat itu juga}, untuk pertanyaan yang belum pernah ada sebelumnya.
\end{iband}
\end{frame}

\framekicker{Dasar}
\begin{frame}[shrink=25]{Penggunanya sering justru bukan programmer}
{\fontsize{9.4}{12.6}\selectfont\color{fg2}Pemakaian agen kode yang paling banyak menambah nilai sering bukan menulis perangkat lunak, melainkan \textbf{menjawab pertanyaan atas data} untuk orang yang tidak menulis kode.\par}\vskip 4pt
\begin{minipage}[t]{0.485\textwidth}
{\fontsize{9.4}{12.6}\selectfont\color{fg2}\textbf{Yang dilihat pengguna}\par}\vskip 4pt
{\fontsize{9.4}{12.6}\selectfont\color{fg2}\textbackslash\{\}u201cBerapa rata-rata tunggakan nasabah kategori B tahun lalu, dipisah per wilayah?\textbackslash\{\}u201d --- lalu sebuah tabel dan grafik.\par}\vskip 4pt

\end{minipage}%
\hfill
\begin{minipage}[t]{0.485\textwidth}
{\fontsize{9.4}{12.6}\selectfont\color{fg2}\textbf{Yang terjadi}\par}\vskip 4pt
{\fontsize{9.4}{12.6}\selectfont\color{fg2}Agen menulis kueri dan beberapa baris pengolahan, menjalankannya, dan menyusun hasilnya. Kodenya tersimpan di jejak.\par}\vskip 4pt

\end{minipage}%
\begin{iband}[signal]
Dan itu yang membuatnya bisa dipertanggungjawabkan: \textbf{kalau angkanya dipertanyakan, kodenya bisa dibaca.} Jawaban langsung dari model tidak punya sifat itu.
\end{iband}
\end{frame}

\coursesection{02}{Alatnya}{Lima jenis, dan garis yang memisahkan dua kelompok.}

\framekicker{Alat}
\begin{frame}[shrink=25]{Lima jenis alat kode, dan garisnya}
\begin{center}
\includegraphics[width=0.92\linewidth,height=0.60\textheight,keepaspectratio]{figs/agents10-tools.tex.pdf}
\end{center}
\vskip 3pt{\color{ink3}\fontsize{7.4}{9.4}\selectfont Membaca aman dan murah; menulis dan menjalankan tidak.\par}
\begin{itable}{>{\raggedright\arraybackslash}p{0.2256\textwidth}>{\raggedright\arraybackslash}p{0.1692\textwidth}>{\raggedright\arraybackslash}p{0.5452\textwidth}}
\thead{Alat} & \thead{Kelompok} & \thead{Catatan} \\ \midrule
Peta repo & Baca & Struktur repo dalam beberapa ribu token \\
Cari kode & Baca & Berdasarkan nama, simbol, atau teks --- bukan kemiripan makna \\
Baca berkas & Baca & Sebagian berkas, dengan nomor baris \\
Sunting berkas & \textbf{Tulis} & Butuh gerbang, dan harus bisa dibatalkan \\
Jalankan perintah & \textbf{Tulis} & Paling berbahaya: ia bisa jadi alat apa pun \\
\end{itable}
\end{frame}

\framekicker{Alat}
\begin{frame}[shrink=25]{Kueri basis data sebagai alat, dan batas yang harus menyertainya}
{\fontsize{9.4}{12.6}\selectfont\color{fg2}Memberi agen kemampuan menulis kueri sangat kuat untuk pertanyaan analitis --- dan membawa risiko yang bentuknya sudah dikenal orang basis data.\par}\vskip 4pt
\begin{minipage}[t]{0.494\textwidth}
\begin{icardgood}
\cardh{Hanya baca, selalu}
Akun terpisah yang benar-benar tidak punya izin menulis. Bukan janji di prompt --- izin di basis datanya.
\end{icardgood}
\end{minipage}%
\hfill
\begin{minipage}[t]{0.494\textwidth}
\begin{icardgood}
\cardh{Batas waktu dan baris}
Kueri tanpa batas bisa menahan basis data produksi. Batasi di sisi peladen, bukan di kueri.
\end{icardgood}
\end{minipage}%


\vskip 4pt

\begin{minipage}[t]{0.494\textwidth}
\begin{icardgood}
\cardh{Tampilan, bukan tabel mentah}
Agen melihat tampilan yang sudah menyaring kolom pribadi. Yang tidak terlihat tidak bisa bocor.
\end{icardgood}
\end{minipage}%
\hfill
\begin{minipage}[t]{0.494\textwidth}
\begin{icardgood}
\cardh{Kuerinya masuk jejak}
Ini bukti terbaik yang bisa Anda punya untuk angka mana pun yang dilaporkan.
\end{icardgood}
\end{minipage}%
\end{frame}

\coursesection{03}{Menemukan yang relevan}{Sebab seluruh repo tidak akan pernah muat.}

\framekicker{Konteks}
\begin{frame}[shrink=25]{Seratus berkas sudah melewati jendela}
\begin{center}
\includegraphics[width=0.97\linewidth,height=0.56\textheight,keepaspectratio]{figs/agents10-repo.tex.pdf}
\end{center}
\vskip 3pt{\color{ink3}\fontsize{7.4}{9.4}\selectfont Seluruh isi berkas lawan peta satu baris per berkas, dihitung. Langkahi menurut ukuran repo.\par}
{\fontsize{9.4}{12.6}\selectfont\color{fg2}\textbackslash\{\}u201cMasukkan saja seluruh kode ke konteks\textbackslash\{\}u201d berhenti bekerja hampir seketika. Petanya dua orde besaran lebih kecil --- dan itulah yang membuat pola \textbf{peta $\rightarrow$ cari $\rightarrow$ baca beberapa} jadi satu-satunya yang bisa diskalakan.\par}\vskip 4pt
\end{frame}

\framekicker{Konteks}
\begin{frame}[shrink=25]{Peta, cari, baca sedikit}
\begin{center}
\includegraphics[width=0.92\linewidth,height=0.60\textheight,keepaspectratio]{figs/agents10-find.tex.pdf}
\end{center}
\vskip 3pt{\color{ink3}\fontsize{7.4}{9.4}\selectfont Tiga langkah, dan hanya langkah terakhir yang mahal.\par}
{\fontsize{9.4}{12.6}\selectfont\color{fg2}Dua langkah pertama murah dan bisa diulang; hanya langkah ketiga yang memakan konteks. Itu sebabnya urutannya penting.\par}\vskip 4pt
\end{frame}

\framekicker{Konteks}
\begin{frame}[shrink=25]{\ldots{}dan apa yang dikerjakan tiap langkah}
\begin{isteps}
\item \{'h': 'Peta memberi tahu apa yang ADA', 'p': 'Nama berkas, kelas, fungsi utama. Cukup untuk memutuskan ke mana harus mencari.'\}
\item \{'h': 'Pencarian menyempitkan', 'p': 'Berdasarkan nama simbol atau teks --- dan di sini pencarian teks biasa \textbf{mengalahkan} pencarian kemiripan makna, sebab nama fungsi adalah nama, bukan konsep.'\}
\item \{'h': 'Baca tiga sampai lima berkas', 'p': 'Itu yang masuk konteks. Sisanya tidak pernah dibaca, dan itu bukan kerugian.'\}
\end{isteps}
\begin{iband}[signal]
Pola yang persis sama dengan Bab 4 dan Bab 9: \textbf{jaring lebar dengan yang murah, saringan halus dengan yang mahal.} Ia muncul untuk ketiga kalinya karena alasan yang sama.
\end{iband}
\end{frame}

\framekicker{Konteks}
\begin{frame}[shrink=25]{Memadatkan konteks tanpa kehilangan yang menentukan}
{\fontsize{9.4}{12.6}\selectfont\color{fg2}Sesi kerja yang panjang akan memenuhi jendela. Bab 4 sudah menyebut bahayanya meringkas berulang; untuk kode ada beberapa hal yang \textbf{tidak boleh} hilang dalam pemadatan.\par}\vskip 4pt
\begin{minipage}[t]{0.494\textwidth}
\begin{icardgood}
\cardh{Harus bertahan}
Tujuan awal, berkas yang sedang disunting, hasil uji terakhir, dan keputusan yang sudah diambil beserta alasannya.
\end{icardgood}
\end{minipage}%
\hfill
\begin{minipage}[t]{0.494\textwidth}
\begin{icard}
\cardh{Boleh hilang}
Isi berkas yang sudah dibaca (bisa dibaca ulang), keluaran perintah yang sudah tidak relevan, jalan buntu yang sudah ditinggalkan.
\end{icard}
\end{minipage}%
\begin{iband}[signal]
Aturan yang membedakan keduanya sama seperti Bab 4: \textbf{apa yang bisa diambil ulang dengan alat boleh dibuang; apa yang hanya ada di percakapan harus dipertahankan.}
\end{iband}
\end{frame}

\framekicker{Konteks}
\begin{frame}[shrink=25]{Menjaga singgahan tetap hidup di sesi kode}
{\fontsize{9.4}{12.6}\selectfont\color{fg2}Bab 2 menyebut singgahan prompt sebagai penghematan terbesar yang tersedia. Pada agen kode ia lebih penting lagi, sebab sesinya panjang dan bagian tetapnya besar --- perintah sistem, skema alat, dan peta repo.\par}\vskip 4pt
\begin{isteps}
\item \{'h': 'Taruh yang tetap di depan, dan jangan diutak-atik', 'p': 'Peta repo yang disusun ulang tiap giliran membunuh singgahan tanpa satu pun pesan galat.'\}
\item \{'h': 'Tambahkan di belakang, jangan menyisipkan di tengah', 'p': 'Menyisipkan sesuatu di tengah mengubah awalan semua yang sesudahnya.'\}
\item \{'h': 'Periksa angkanya, jangan percaya', 'p': '\inlinecode{cache\_read\_input\_tokens} harus bukan nol. Kalau jadi nol, sesuatu di awalan berubah.'\}
\end{isteps}
\end{frame}

\framekicker{Alat}
\begin{frame}[shrink=25]{Menyunting: tambalan, bukan tulis ulang}
{\fontsize{9.4}{12.6}\selectfont\color{fg2}Ada dua cara alat suntingan bisa dirancang, dan pilihannya mengubah banyak hal sekaligus.\par}\vskip 4pt
\begin{itable}{>{\raggedright\arraybackslash}p{0.2632\textwidth}>{\raggedright\arraybackslash}p{0.2444\textwidth}>{\raggedright\arraybackslash}p{0.4324\textwidth}}
\thead{Cara} & \thead{Biaya token} & \thead{Risikonya} \\ \midrule
Kembalikan seluruh berkas & Besar --- dua kali isi berkas & Bagian yang tidak disentuh bisa ikut berubah tanpa disadari \\
\textbf{Tambalan berupa diff} & Kecil --- hanya yang berubah & Bisa gagal diterapkan kalau berkasnya bergeser; itu \textbf{kegagalan yang kelihatan} \\
\end{itable}
\begin{iband}[signal]
Baris kedua lebih baik karena alasan yang tidak langsung terlihat: \textbf{diff yang gagal diterapkan adalah galat}, sedangkan tulis ulang yang diam-diam mengubah baris lain adalah bug yang lolos ke tinjauan.
\end{iband}
\end{frame}

\framekicker{Alat}
\begin{frame}[shrink=25]{Mencari kode: teks mengalahkan makna}
{\fontsize{9.4}{12.6}\selectfont\color{fg2}Bab 4 menyebut kelemahan pencarian kemiripan makna untuk angka dan pengenal. Pada kode, kelemahan itu jadi menentukan: \textbf{nama fungsi adalah nama, bukan konsep.}\par}\vskip 4pt
\begin{minipage}[t]{0.3253\textwidth}
\begin{icardgood}
\cardh{Cari teks persis}
Nama simbol, pemanggilan, string galat. Cepat, pasti, dan hampir selalu yang benar.
\end{icardgood}
\end{minipage}%
\hfill
\begin{minipage}[t]{0.3253\textwidth}
\begin{icardgood}
\cardh{Cari menurut struktur}
Definisi, pemanggil, penurunan kelas. Lebih tepat lagi kalau bahasanya mendukung.
\end{icardgood}
\end{minipage}%
\hfill
\begin{minipage}[t]{0.3253\textwidth}
\begin{icard}
\cardh{Cari kemiripan makna}
Berguna untuk ``di mana logika kredit?'' dan buruk untuk ``di mana \inlinecode{hitung\_dscr} dipanggil?''
\end{icard}
\end{minipage}%
\begin{iband}[signal]
Bentuk yang paling sering benar: \textbf{cari teks dulu; kemiripan makna hanya kalau pencarian teks tidak menemukan apa-apa} --- dan sebutkan dalam deskripsi alatnya kapan masing-masing dipakai.
\end{iband}
\end{frame}

\framekicker{Alat}
\begin{frame}[shrink=25]{Membaca berkas: sebagian, dengan nomor baris}
\begin{isteps}
\item \{'h': 'Kembalikan potongan, bukan berkas penuh', 'p': 'Berkas 2 000 baris masuk konteks satu kali dan dibayar berkali kali. Sebagian besar tugas butuh tiga puluh baris.'\}
\item \{'h': 'Sertakan nomor barisnya', 'p': 'Supaya tambalan bisa merujuk lokasi, dan supaya jejaknya bisa dibaca ulang.'\}
\item \{'h': 'Sebutkan kalau dipotong', 'p': 'Berkas terpotong yang menyamar sebagai berkas penuh menghasilkan kesimpulan yang salah tentang apa yang ada di dalamnya.'\}
\end{isteps}
\begin{iband}[signal]
Ketiganya pengulangan aturan Bab 5 tentang mengolah keluaran alat, diterapkan pada bentuk data yang paling sering membanjiri konteks agen kode.
\end{iband}
\end{frame}

\coursesection{04}{Dinding}{Satu-satunya alat di modul ini yang butuh isolasi sungguhan.}

\framekicker{Sandbox}
\begin{frame}[shrink=25]{Alat yang bisa menjadi alat apa pun}
{\fontsize{9.4}{12.6}\selectfont\color{fg2}Semua alat lain di modul ini punya batas yang jelas: \inlinecode{ambil\_data} mengambil data. Alat yang menjalankan kode tidak punya batas seperti itu --- \textbf{ia bisa melakukan apa pun yang bisa dilakukan kode} di lingkungan tempat ia berjalan.\par}\vskip 4pt
\begin{iband}[signal]
Karena itu pertanyaan keamanannya berubah bentuk. Untuk alat lain: \textit{apa yang boleh dipanggil?} Untuk alat kode: \textbf{di mana ia berjalan, dan apa yang bisa dijangkau dari sana?}
\end{iband}
{\fontsize{9.4}{12.6}\selectfont\color{fg2}Dan itu pertanyaan infrastruktur, bukan pertanyaan prompt --- dijawab dengan proses terpisah, jaringan yang ditutup, dan lingkungan yang tidak memuat kredensial apa pun.\par}\vskip 4pt
\end{frame}

\framekicker{Sandbox}
\begin{frame}[shrink=25]{Berlapis, dan jujur soal yang tidak ditutupnya}
\begin{itable}{>{\raggedright\arraybackslash}p{0.2256\textwidth}>{\raggedright\arraybackslash}p{0.3384\textwidth}>{\raggedright\arraybackslash}p{0.376\textwidth}}
\thead{Lapis} & \thead{Menahan} & \thead{TIDAK menahan} \\ \midrule
Tanpa kredensial di lingkungan & Pemakaian kunci yang bocor & Kunci yang ada di dalam kode yang dibacanya \\
Jaringan ditutup & Pengiriman data keluar & Apa pun yang sudah ada di dalam sandbox \\
Batas waktu dan memori & Gelung tak berujung & Kerusakan yang selesai dalam satu detik \\
Proses / kontainer terpisah & Sebagian besar hal & Celah pada kernel atau salah konfigurasi \\
Berkas hanya-baca & Perubahan yang tidak disengaja & Pembacaan berkas yang seharusnya tidak terbaca \\
\end{itable}
\begin{iband}[signal]
Kolom ketiga yang membuat tabel ini berguna. \textbf{Sandbox yang dijelaskan sebagai \textbackslash\{\}u201caman\textbackslash\{\}u201d membuat orang berhenti berpikir}; yang menyebut batasnya membuat orang menaruh data sensitif di tempat lain.
\end{iband}
\end{frame}

\framekicker{Sandbox}
\begin{frame}[shrink=25]{Menjalankan perintah bukan hal yang sama dengan menjalankan kode}
\begin{minipage}[t]{0.485\textwidth}
{\fontsize{9.4}{12.6}\selectfont\color{fg2}\textbf{Penerjemah kode}\par}\vskip 4pt
{\fontsize{9.4}{12.6}\selectfont\color{fg2}Satu bahasa, pustaka yang bisa dibatasi, keluaran yang bentuknya diketahui. Bisa dikurung relatif rapat.\par}\vskip 4pt

\end{minipage}%
\hfill
\begin{minipage}[t]{0.485\textwidth}
{\fontsize{9.4}{12.6}\selectfont\color{fg2}\textbf{Baris perintah}\par}\vskip 4pt
{\fontsize{9.4}{12.6}\selectfont\color{fg2}Seluruh sistem operasi, termasuk hal yang tidak terpikir saat merancang. Daftar perintah yang diizinkan lebih masuk akal daripada daftar yang dilarang.\par}\vskip 4pt

\end{minipage}%
\begin{iband}[signal]
Kalau harus memilih satu untuk dimulai, mulai dari \textbf{penerjemah kode tanpa jaringan}. Ia mencakup sebagian besar kegunaan analitis dengan permukaan yang jauh lebih kecil.
\end{iband}
\end{frame}

\coursesection{05}{Untuk rekayasa perangkat lunak}{Di sini pemeriksanya sudah ada, dan itu mengubah segalanya.}

\framekicker{Rekayasa}
\begin{frame}[shrink=25]{Gelung yang punya pemeriksa objektif}
\begin{center}
\includegraphics[width=0.92\linewidth,height=0.60\textheight,keepaspectratio]{figs/agents10-loop.tex.pdf}
\end{center}
\vskip 3pt{\color{ink3}\fontsize{7.4}{9.4}\selectfont Sunting, jalankan uji, ulangi --- dan ujinya yang menentukan berhenti.\par}
{\fontsize{9.4}{12.6}\selectfont\color{fg2}Bab 3 dan Bab 6 menyebut syarat refleksi bekerja: harus ada pemeriksa yang lebih mudah daripada menjawab. Pemrograman adalah domain yang \textbf{sudah punya itu sejak awal} --- uji yang dijalankan, penyusun yang mengeluh, linter yang menolak.\par}\vskip 4pt
\begin{iband}[signal]
Itu sebabnya agen kode terasa jauh lebih berhasil daripada agen di domain lain, dan \textbf{bukan karena modelnya lebih pandai soal kode}. Ia bekerja di tempat yang jawabannya bisa diperiksa mesin.
\end{iband}
\end{frame}

\framekicker{Rekayasa}
\begin{frame}[shrink=25]{Alur tetap sering mengalahkan agen bebas}
\begin{center}

\includegraphics[width=0.92\linewidth,height=0.60\textheight,keepaspectratio]{figs/agents10-agentless.tex.pdf}
\end{center}
\vskip 3pt{\color{ink3}\fontsize{7.4}{9.4}\selectfont Empat langkah tetap, lawan gelung yang memilih sendiri.\par}
{\fontsize{9.4}{12.6}\selectfont\color{fg2}Untuk perbaikan bug yang bentuknya berulang --- temukan berkas, usulkan beberapa tambalan, jalankan uji, ambil yang lulus --- \textbf{alur tetap sering mengungguli agen bebas}, dengan biaya lebih kecil dan hasil yang bisa diramalkan.\par}\vskip 4pt
\begin{iband}[signal]
Ini pengulangan pesan Bab 1 pada domain yang paling banyak menghasilkan demo agen: \hilite{kalau langkahnya sudah diketahui, yang Anda butuhkan alur tetap}, dan itu tetap berlaku di sini.
\end{iband}
\end{frame}

\framekicker{Rekayasa}
\begin{frame}[shrink=25]{Membaca angka tolok ukur kode dengan curiga}
\begin{minipage}[t]{0.494\textwidth}
\begin{icardbad}
\cardh{Kontaminasi}
Repo publik dan isunya kemungkinan besar ada di data latih. Angka tinggi pada isu lama tidak berarti banyak.
\end{icardbad}
\end{minipage}%
\hfill
\begin{minipage}[t]{0.494\textwidth}
\begin{icardbad}
\cardh{Ujinya yang menilai}
\textbackslash\{\}u201cBerhasil\textbackslash\{\}u201d berarti uji yang ada lulus --- bukan bahwa perbaikannya benar. Tambalan yang mematikan ujinya juga lulus.
\end{icardbad}
\end{minipage}%


\vskip 4pt

\begin{minipage}[t]{0.494\textwidth}
\begin{icardbad}
\cardh{Berapa percobaan?}
pass@k dengan k besar (Bab 7) membuat angka melonjak tanpa keandalan bertambah.
\end{icardbad}
\end{minipage}%
\hfill
\begin{minipage}[t]{0.494\textwidth}
\begin{icardbad}
\cardh{Bukan repo Anda}
Basis kode besar yang lama, tanpa uji, dengan konvensi sendiri adalah masalah yang berbeda.
\end{icardbad}
\end{minipage}%
\end{frame}

\framekicker{Rekayasa}
\begin{frame}[shrink=25]{Basis kode yang sudah ada adalah masalah yang berbeda}
\begin{minipage}[t]{0.485\textwidth}
{\fontsize{9.4}{12.6}\selectfont\color{fg2}\textbf{Proyek baru}\par}\vskip 4pt
{\fontsize{9.4}{12.6}\selectfont\color{fg2}Tidak ada konvensi yang harus dipatuhi, tidak ada kode lama yang bisa rusak, dan uji ditulis bersamaan. Di sinilah demo terlihat mengesankan.\par}\vskip 4pt

\end{minipage}%
\hfill
\begin{minipage}[t]{0.485\textwidth}
{\fontsize{9.4}{12.6}\selectfont\color{fg2}\textbf{Basis kode yang berumur}\par}\vskip 4pt
{\fontsize{9.4}{12.6}\selectfont\color{fg2}Konvensi tak tertulis, kaitan yang tidak terlihat, uji yang menutup sebagian kecil saja. Di sinilah pekerjaan sebenarnya berada.\par}\vskip 4pt

\end{minipage}%
{\fontsize{9.4}{12.6}\selectfont\color{fg2}Yang paling menolong pada kolom kanan bukan model yang lebih pandai, melainkan \textbf{konteks yang lebih baik}: berkas panduan konvensi, peta modul, dan contoh perubahan serupa yang pernah diterima.\par}\vskip 4pt
\begin{iband}[signal]
Pola yang sama dengan seluruh modul: \textbf{perbaiki apa yang masuk ke konteks sebelum mengganti modelnya.}
\end{iband}
\end{frame}

\framekicker{Rekayasa}
\begin{frame}[shrink=25]{Perubahan yang menyentuh banyak berkas}
{\fontsize{9.4}{12.6}\selectfont\color{fg2}Perubahan satu berkas relatif mudah. Perubahan yang harus konsisten di sepuluh berkas adalah tempat agen paling sering gagal separuh jalan --- dan separuh jalan pada kode berarti \textbf{repo yang tidak bisa dibangun}.\par}\vskip 4pt
\begin{isteps}
\item \{'h': 'Kumpulkan seluruh tambalan dulu, terapkan di akhir', 'p': 'Sama seperti Bab 8: tunda tulisan sampai semuanya siap, supaya kegagalan sebagian tidak meninggalkan keadaan setengah jadi.'\}
\item \{'h': 'Jalankan pembangunan dan uji sesudah semuanya diterapkan', 'p': 'Bukan sesudah tiap berkas --- itu memberi sinyal yang menyesatkan.'\}
\item \{'h': 'Kalau gagal, kembalikan semuanya', 'p': 'Cabang terpisah membuat ini satu perintah, bukan satu penyelamatan.'\}
\end{isteps}
\end{frame}

\framekicker{Rekayasa}
\begin{frame}[shrink=25]{Meninjau kode yang ditulis agen}
\begin{minipage}[t]{0.494\textwidth}
\begin{icardaccent}
\cardh{Baca diff-nya, bukan ringkasannya}
Ringkasan yang ditulis agen menjelaskan apa yang \textbf{dimaksudkannya}. Diff menunjukkan apa yang \textbf{dilakukannya}.
\end{icardaccent}
\end{minipage}%
\hfill
\begin{minipage}[t]{0.494\textwidth}
\begin{icardaccent}
\cardh{Periksa apakah ujinya ikut berubah}
Uji yang disesuaikan supaya lulus adalah temuan, bukan detail.
\end{icardaccent}
\end{minipage}%


\vskip 4pt

\begin{minipage}[t]{0.494\textwidth}
\begin{icard}
\cardh{Perhatikan yang DIHAPUS}
Penanganan galat dan pemeriksaan yang hilang jarang terlihat di ringkasan mana pun.
\end{icard}
\end{minipage}%
\hfill
\begin{minipage}[t]{0.494\textwidth}
\begin{icard}
\cardh{Tolak perubahan yang terlalu besar}
Diff yang tidak bisa dibaca dalam sepuluh menit tidak akan ditinjau dengan sungguh-sungguh --- oleh siapa pun.
\end{icard}
\end{minipage}%
\begin{iband}[signal]
Kartu terakhir sering menentukan apakah agen kode benar-benar menghemat waktu: \textbf{perubahan kecil yang sering mengalahkan perubahan besar yang jarang}, sebab hanya yang pertama benar-benar ditinjau.
\end{iband}
\end{frame}

\framekicker{Rekayasa}
\begin{frame}[shrink=25]{Merencanakan sebelum menyunting}
{\fontsize{9.4}{12.6}\selectfont\color{fg2}Bab 6 menyebut rencana berguna kalau tiap langkahnya bisa dipetakan ke alat dan ditandai selesai. Pada kode, syarat itu mudah dipenuhi --- dan rencananya bisa ditunjukkan sebelum ada perubahan apa pun.\par}\vskip 4pt
\begin{isteps}
\item \{'h': 'Sebutkan berkas yang akan disentuh', 'p': 'Sebelum menyentuhnya. Peninjau bisa menghentikan arah yang salah ketika biayanya masih nol.'\}
\item \{'h': 'Sebutkan uji yang harus lulus', 'p': 'Kalau tidak ada ujinya, langkah pertama adalah menulis ujinya.'\}
\item \{'h': 'Sebutkan apa yang TIDAK akan diubah', 'p': 'Ini yang paling menenangkan peninjau, dan paling jarang dituliskan.'\}
\end{isteps}
\begin{iband}[signal]
Bentuk pengawasan yang paling bisa dijalankan tetap sama seperti Bab 6: \textbf{satu tinjauan rencana yang benar-benar dibaca, bukan sepuluh persetujuan yang diklik.}
\end{iband}
\end{frame}

\coursesection{06}{Melatih model kode}{Satu gelung yang berjalan karena pemeriksanya gratis.}

\framekicker{Melatih}
\begin{frame}[shrink=25]{Data buatan yang bisa diperiksa sendiri}
{\fontsize{9.4}{12.6}\selectfont\color{fg2}Kode punya sifat yang jarang: \textbf{benar atau tidaknya bisa diperiksa mesin, gratis, dan berkali-kali.} Itu membuat gelung yang tidak mungkin di domain lain jadi mungkin di sini.\par}\vskip 4pt
\begin{isteps}
\item \{'h': 'Bangkitkan soal dan calon jawabannya', 'p': 'Dari repo yang ada, atau dari soal yang disusun model sendiri.'\}
\item \{'h': 'Jalankan ujinya', 'p': 'Yang lulus disimpan, yang gagal dibuang. Tidak ada manusia yang perlu menilai.'\}
\item \{'h': 'Latih pada yang lulus, ulangi', 'p': 'Model membaik, soal yang bisa diselesaikannya bertambah, dan putaran berikutnya menghasilkan data yang lebih baik.'\}
\end{isteps}
\begin{iband}[signal]
Bab 3 menyebut batas gelung semacam ini, dan batas itu berlaku di sini juga: \textbf{ia hanya sekuat pemeriksanya}, dan soal yang dibangkitkan cenderung mirip yang sudah bisa dijawab.
\end{iband}
\end{frame}

\framekicker{Melatih}
\begin{frame}[shrink=25]{Kenapa gelung itu berhenti bekerja di luar kode}
{\fontsize{9.4}{12.6}\selectfont\color{fg2}Gelung data buatan tadi terlihat seperti resep umum, dan ia bukan. Ia bergantung pada satu sifat yang jarang: \textbf{pemeriksa yang gratis, otomatis, dan tidak bisa dibujuk.}\par}\vskip 4pt
\begin{minipage}[t]{0.485\textwidth}
{\fontsize{9.4}{12.6}\selectfont\color{fg2}\textbf{Punya pemeriksa seperti itu}\par}\vskip 4pt
\begin{ibullets}
\item Kode --- uji dijalankan
\item Matematika --- hasilnya dicocokkan
\item Kueri --- dibandingkan dengan jawaban yang diketahui
\end{ibullets}

\end{minipage}%
\hfill
\begin{minipage}[t]{0.485\textwidth}
{\fontsize{9.4}{12.6}\selectfont\color{fg2}\textbf{Tidak punya}\par}\vskip 4pt
\begin{ibullets}
\item Menilai kelayakan kredit
\item Meringkas dokumen
\item Menjawab pertanyaan kebijakan
\end{ibullets}

\end{minipage}%
\begin{iband}[signal]
Untuk kolom kanan, yang menggantikan pemeriksa otomatis adalah \textbf{kumpulan uji yang ditulis manusia} --- Bab 7 --- dan tidak ada jalan pintas yang menghindarinya.
\end{iband}
\end{frame}

\framekicker{Risiko}
\begin{frame}[shrink=25]{Kode yang berjalan bukan kode yang benar}
{\fontsize{9.4}{12.6}\selectfont\color{fg2}Gelung sunting-uji berhenti ketika ujinya hijau. Itu syarat yang jelas dan \textbf{bukan} definisi kebenaran --- ia hanya sekuat cakupan ujinya.\par}\vskip 4pt
\begin{minipage}[t]{0.3253\textwidth}
\begin{icardbad}
\cardh{Yang tidak diuji}
Perilaku yang tidak punya uji bisa rusak tanpa satu pun lampu merah menyala.
\end{icardbad}
\end{minipage}%
\hfill
\begin{minipage}[t]{0.3253\textwidth}
\begin{icardbad}
\cardh{Kasus tepi}
Uji yang ada biasanya menutup jalur bahagia. Kasus tepi justru yang paling sering rusak.
\end{icardbad}
\end{minipage}%
\hfill
\begin{minipage}[t]{0.3253\textwidth}
\begin{icardbad}
\cardh{Sifat non-fungsional}
Kinerja, penggunaan memori, dan keamanan hampir tidak pernah punya uji.
\end{icardbad}
\end{minipage}%
\begin{iband}[signal]
Karena itu tinjauan manusia tetap wajib, dan pertanyaannya bukan ``apakah ujinya lulus'' melainkan \textbf{``apa yang bisa rusak dan tidak akan ketahuan''}.
\end{iband}
\end{frame}

\framekicker{Praktik}
\begin{frame}[shrink=25]{Menyimpan konteks proyek supaya tidak diulang tiap sesi}
{\fontsize{9.4}{12.6}\selectfont\color{fg2}Hal yang sama dijelaskan berulang tiap sesi --- konvensi, cara menjalankan uji, bagian mana yang rapuh --- adalah biaya yang bisa dihapus sekali.\par}\vskip 4pt
\begin{isteps}
\item \{'h': 'Satu berkas panduan di repo', 'p': 'Konvensi, perintah yang dipakai, dan hal yang tidak boleh disentuh. Ditinjau seperti kode, diversikan seperti kode.'\}
\item \{'h': 'Taruh di bagian tetap konteks', 'p': 'Supaya ikut disinggahkan, bukan dikirim ulang sebagai pesan baru tiap kali.'\}
\item \{'h': 'Perbarui ketika agen salah menebak', 'p': 'Tiap kesalahan yang berasal dari konvensi yang tidak tertulis adalah satu baris yang kurang di berkas itu.'\}
\end{isteps}
\begin{iband}[signal]
Langkah ketiga yang membuatnya membaik sendiri: \textbf{berkas panduan yang tumbuh dari kesalahan nyata} jauh lebih berguna daripada yang ditulis sekaligus di awal.
\end{iband}
\end{frame}

\coursesection{07}{Risiko dan praktik}{Yang khas agen kode, dan tidak muncul pada agen lain.}

\framekicker{Risiko}
\begin{frame}[shrink=25]{Empat risiko yang khas agen kode}
\begin{minipage}[t]{0.494\textwidth}
\begin{icardbad}
\cardh{Rahasia di dalam kode}
Agen yang membaca repo membaca juga kunci yang tertinggal di sana --- dan bisa memasukkannya ke konteks, lalu ke jejak.
\end{icardbad}
\end{minipage}%
\hfill
\begin{minipage}[t]{0.494\textwidth}
\begin{icardbad}
\cardh{Ketergantungan yang ditambahkannya}
\inlinecode{pip install} sesuatu yang namanya mirip pustaka asli adalah kelas serangan tersendiri.
\end{icardbad}
\end{minipage}%


\vskip 4pt

\begin{minipage}[t]{0.494\textwidth}
\begin{icardbad}
\cardh{Menyesuaikan uji, bukan kode}
Cara tercepat membuat uji hijau adalah mengubah ujinya. Ini terjadi, dan tertangkap hanya kalau diff-nya dibaca.
\end{icardbad}
\end{minipage}%
\hfill
\begin{minipage}[t]{0.494\textwidth}
\begin{icardbad}
\cardh{Perubahan yang merusak diam-diam}
Menghapus penanganan galat supaya lulus uji terlihat seperti penyederhanaan.
\end{icardbad}
\end{minipage}%
{\fontsize{9.4}{12.6}\selectfont\color{fg2}Keempatnya punya satu penawar yang sama dan membosankan: \textbf{setiap perubahan lewat tinjauan manusia, sebagai diff yang dibaca} --- persis seperti kode dari orang lain.\par}\vskip 4pt
\end{frame}

\framekicker{Risiko}
\begin{frame}[shrink=25]{Rahasia: cegah di dua tempat, bukan satu}
\begin{isteps}
\item \{'h': 'Jangan ada rahasia di repo', 'p': 'Ini seharusnya sudah benar sebelum ada agen, dan agen membuat akibatnya lebih cepat terasa.'\}
\item \{'h': 'Saring saat masuk konteks', 'p': 'Pola kunci yang dikenali dibuang sebelum masuk --- sebab yang masuk konteks akan masuk jejak.'\}
\item \{'h': 'Saring saat masuk jejak', 'p': 'Jejak disimpan lama dan dibaca banyak orang. Ini tempat kebocoran yang paling sering tidak terpikir.'\}
\end{isteps}
\begin{iband}[signal]
Perhatikan langkah ketiga: \textbf{jejak yang dirancang untuk audit bisa berubah jadi tempat penyimpanan rahasia} kalau tidak ada yang menyaringnya.
\end{iband}
\end{frame}

\framekicker{Praktik}
\begin{frame}[shrink=25]{Urutan memasang agen kode dengan aman}
\begin{isteps}
\item \{'h': 'Mulai dari baca saja', 'p': 'Peta, cari, baca. Agen yang menjelaskan basis kode sudah berguna dan tidak bisa merusak apa pun.'\}
\item \{'h': 'Tambahkan penerjemah kode tanpa jaringan', 'p': 'Ini membuka sebagian besar kegunaan analitis dengan permukaan terkecil.'\}
\item \{'h': 'Baru sunting berkas, dengan diff yang ditinjau', 'p': 'Dan tidak pernah langsung ke cabang utama.'\}
\item \{'h': 'Baris perintah paling akhir, dengan daftar yang diizinkan', 'p': 'Kalau memang diperlukan. Banyak sistem berhenti sebelum ini dan sudah cukup.'\}
\end{isteps}
\begin{iband}[signal]
Urutan ini sama bentuknya dengan Bab 5: \textbf{alat baca dulu, alat tulis satu per satu dengan alasan tertulis.} Yang berubah cuma seberapa tajam alatnya.
\end{iband}
\end{frame}

\framekicker{Praktik}
\begin{frame}[shrink=25]{Yang ada di demo, dan apa yang jujur disebutkan tentangnya}
{\fontsize{9.4}{12.6}\selectfont\color{fg2}Repo demo punya satu kasus yang menjalankan kode di balik sandbox berlapis --- dan dokumennya menyebut \textbf{apa yang tidak ditutup lapisan-lapisan itu}, bukan menyatakannya aman.\par}\vskip 4pt
{\fontsize{9.4}{12.6}\selectfont\color{fg2}Itu keputusan penulisan yang disengaja, dan pantas ditiru: pembaca yang diberi tahu batasnya akan menaruh data sensitif di tempat lain. Pembaca yang diberi tahu \textbackslash\{\}u201caman\textbackslash\{\}u201d tidak akan.\par}\vskip 4pt
\begin{iband}[signal]
Sama seperti aplikasi selulernya, yang dokumennya menyebut versi Flutter yang dipakai untuk membangunnya dan apa yang belum diuji. \textbf{Klaim yang bisa diperiksa mengalahkan klaim yang menenangkan.}
\end{iband}
\end{frame}

\framekicker{Praktik}
\begin{frame}[shrink=25]{Berapa biayanya, dan di mana biayanya}
{\fontsize{9.4}{12.6}\selectfont\color{fg2}Agen kode adalah beban konteks yang berat: peta repo, berkas yang dibaca, keluaran uji, dan diff --- semuanya masuk riwayat dan dibayar ulang tiap giliran.\par}\vskip 4pt
\begin{itable}{>{\raggedright\arraybackslash}p{0.2444\textwidth}>{\raggedright\arraybackslash}p{0.3196\textwidth}>{\raggedright\arraybackslash}p{0.376\textwidth}}
\thead{Sumber biaya} & \thead{Besar kalau} & \thead{Obatnya} \\ \midrule
Berkas yang dibaca & Dibaca penuh, bukan sebagian & Potongan bernomor baris \\
Keluaran uji & Seluruh keluaran masuk, termasuk yang lulus & Kembalikan yang gagal saja \\
Peta repo & Repo besar & Peta per direktori, dimuat saat diperlukan \\
Iterasi & Gelung sunting-uji berputar banyak & Anggaran putaran, dan uji yang lebih cepat \\
\end{itable}
\begin{iband}[signal]
Baris kedua yang paling mudah diperbaiki dan paling sering terlewat: \textbf{keluaran uji yang lulus tidak memberi informasi apa pun}, dan bisa ribuan token.
\end{iband}
\end{frame}

\framekicker{Praktik}
\begin{frame}[shrink=25]{Angka yang memberi tahu agen kode Anda sehat}
\begin{itable}{>{\raggedright\arraybackslash}p{0.2632\textwidth}>{\raggedright\arraybackslash}p{0.282\textwidth}>{\raggedright\arraybackslash}p{0.3948\textwidth}}
\thead{Angka} & \thead{Sehat kalau} & \thead{Gejala kalau tidak} \\ \midrule
Diff yang diterima tanpa perubahan & Naik pelan & Turun $\rightarrow$ konteksnya memburuk, bukan modelnya \\
Putaran sunting-uji per tugas & Kecil dan stabil & Naik $\rightarrow$ ujinya lambat, atau galatnya tidak informatif \\
Berkas dibaca per tugas & Sedikit & Banyak $\rightarrow$ pencariannya tidak menemukan yang tepat \\
Uji yang ikut berubah & \textbf{Nol} & Bukan nol $\rightarrow$ periksa satu per satu, ini bukan detail \\
\end{itable}
{\fontsize{9.4}{12.6}\selectfont\color{fg2}Baris terakhir pantas jadi pemeriksaan otomatis: \textbf{tandai tiap diff yang menyentuh berkas uji}, dan minta alasan tertulis. Sering sah; kadang tidak.\par}\vskip 4pt
\end{frame}

\framekicker{Praktik}
\begin{frame}[shrink=25]{Kapan agen kode adalah alat yang salah}
\begin{minipage}[t]{0.494\textwidth}
\begin{icardbad}
\cardh{Keputusan arsitektur}
Pilihan yang akibatnya bertahun-tahun dan pertukarannya tidak bisa diuji. Agen bisa menuliskan pilihannya; ia tidak bisa menanggungnya.
\end{icardbad}
\end{minipage}%
\hfill
\begin{minipage}[t]{0.494\textwidth}
\begin{icardbad}
\cardh{Kode yang tidak dipahami siapa pun}
Kalau tidak ada manusia yang bisa meninjau hasilnya, kecepatannya tidak berarti apa-apa.
\end{icardbad}
\end{minipage}%


\vskip 4pt

\begin{minipage}[t]{0.494\textwidth}
\begin{icardbad}
\cardh{Perbaikan darurat di produksi}
Tekanan waktu adalah keadaan terburuk untuk meninjau diff dengan sungguh-sungguh.
\end{icardbad}
\end{minipage}%
\hfill
\begin{minipage}[t]{0.494\textwidth}
\begin{icardgood}
\cardh{Di sinilah ia menang}
Perubahan berulang, migrasi mekanis, menulis uji, menjelaskan kode yang sudah ada, dan analisis data.
\end{icardgood}
\end{minipage}%
\end{frame}

\framekicker{Penutup modul}
\begin{frame}[shrink=25]{Sepuluh bab, satu kalimat yang paling sering berlaku}
{\fontsize{9.4}{12.6}\selectfont\color{fg2}Kalau hanya satu hal yang dibawa dari sepuluh bab ini: \textbf{batas kemampuan sebuah agen ada di alat yang dimilikinya} --- dan itu satu-satunya hal tentang sistem berbasis model yang jawabannya berupa daftar, bukan statistik.\par}\vskip 4pt
\begin{isteps}
\item \{'h': 'Karena itu bisa diuji', 'p': 'Satu uji yang gagal kalau seseorang menambahkan alat yang melanggar batasnya.'\}
\item \{'h': 'Karena itu bisa dijelaskan', 'p': 'Kepada pemeriksa, kepada atasan, kepada orang yang menanggung akibatnya.'\}
\item \{'h': 'Karena itu bertahan', 'p': 'Prompt bisa diubah siapa saja; alat yang tidak ada tetap tidak ada.'\}
\end{isteps}
\begin{iband}[signal]
Sisanya --- memori, rencana, refleksi, banyak agen, multi-modal --- semuanya menambah kemampuan. \textbf{Tidak satu pun mengubah kalimat di atas.}
\end{iband}
\end{frame}

\framekicker{Penutup}
\begin{frame}[shrink=25]{Yang dibawa pulang dari bab ini}
\begin{isteps}
\item \{'h': 'Kode mengubah jawaban dari ditebak jadi dihitung', 'p': 'Dan cara menghitungnya tersimpan untuk diperiksa.'\}
\item \{'h': 'Seluruh repo tidak akan muat; petanya muat', 'p': 'Peta $\rightarrow$ cari $\rightarrow$ baca beberapa. Pola yang sama seperti Bab 4 dan 9.'\}
\item \{'h': 'Ini satu-satunya alat yang butuh dinding sungguhan', 'p': 'Dan dindingnya harus menyebutkan apa yang tidak ditutupnya.'\}
\item \{'h': 'Ia berhasil karena pemeriksanya sudah ada', 'p': 'Bukan karena modelnya lebih pandai soal kode.'\}
\item \{'h': 'Alur tetap masih sering menang', 'p': 'Pesan Bab 1, di domain yang paling banyak menghasilkan demo agen.'\}
\end{isteps}
\end{frame}
\note{Penutup modul: dari sepuluh bab ini, satu kalimat yang paling sering berlaku --- batas kemampuan sebuah agen ada di alat yang dimilikinya, dan itu satu-satunya hal tentang sistem berbasis model yang jawabannya bisa berupa daftar.}

\end{document}
